
%%%%%%%%%%%%%%%%%%%%%%%%%%%%%%%%%%%%%%%%%%%%%%%%%%%%%%%%%%%%%%%%%%%%%%%%%%%%%%%%
%% Preambles
%%%%%%%%%%%%%%%%%%%%%%%%%%%%%%%%%%%%%%%%%%%%%%%%%%%%%%%%%%%%%%%%%%%%%%%%%%%%%%%%

\documentclass[a4paper, 11pt, DIV11, BCOR5mm]{scrartcl}
\usepackage[ngerman]{babel}
\usepackage[utf8]{inputenc}
\usepackage{siunitx}
\usepackage{fancyhdr}
\usepackage{float}
\usepackage{graphicx}
\usepackage[dvipsnames]{xcolor}
\usepackage{hyperref}
\usepackage{ragged2e}
\usepackage{titlesec}
\usepackage{listings}
\usepackage{framed}
\usepackage[linguistics]{forest}
%%

%%%%%%%%%%%%%%%%%%%%%%%%%%%%%%%%%%%%%%%%%%%%%%%%%%%%%%%%%%%%%%%%%%%%%%%%%%%%%%%%
%% General configurations
%%%%%%%%%%%%%%%%%%%%%%%%%%%%%%%%%%%%%%%%%%%%%%%%%%%%%%%%%%%%%%%%%%%%%%%%%%%%%%%%
\sisetup{locale=DE}
\graphicspath{ {Bilder/} }
\setlength{\parskip}{12 pt}
\titleformat*{\section}{\LARGE\bfseries}
\titleformat*{\subsection}{\large\bfseries}

%%%%%%%%%%%%%%%%%%%%%%%%%%%%%%%%%%%%%%%%%%%%%%%%%%%%%%%%%%%%%%%%%%%%%%%%%%%%%%%%
%% Counter definitions
%%%%%%%%%%%%%%%%%%%%%%%%%%%%%%%%%%%%%%%%%%%%%%%%%%%%%%%%%%%%%%%%%%%%%%%%%%%%%%%%
\newcounter{imagecounter}

%%%%%%%%%%%%%%%%%%%%%%%%%%%%%%%%%%%%%%%%%%%%%%%%%%%%%%%%%%%%%%%%%%%%%%%%%%%%%%%%
%% New commands definitions
%%%%%%%%%%%%%%%%%%%%%%%%%%%%%%%%%%%%%%%%%%%%%%%%%%%%%%%%%%%%%%%%%%%%%%%%%%%%%%%%
\newcommand{\bfh}{\textsc{Berner Fachhochschule}}
\newcommand{\modul}[1]{\textsc{#1}}
\newcommand{\thema}{Bitcoin}
\newcommand{\code}[1]{{\ttfamily #1}}
\newcommand{\fdate}[1]{\emph{#1\newline\newline}}

\newcommand{\query}[2]{
    \code{?- #1}\\
    \indent
    \code{#2}
}

\NewDocumentCommand\semester{O{Herbstsemester}m}{#1 #2}

%%%%%%%%%%%%%%%%%%%%%%%%%%%%%%%%%%%%%%%%%%%%%%%%%%%%%%%%%%%%%%%%%%%%%%%%%%%%%%%%
%% Color definitions
%%%%%%%%%%%%%%%%%%%%%%%%%%%%%%%%%%%%%%%%%%%%%%%%%%%%%%%%%%%%%%%%%%%%%%%%%%%%%%%%
\definecolor{mygray}{rgb}{0.5,0.5,0.5}
\definecolor{mymauve}{rgb}{0.58,0,0.82}
\definecolor{myred}{HTML}{C70039}
\definecolor{mygreen}{HTML}{C5D200}%{A9C52F}
\definecolor{mydarkgreen}{HTML}{51710A}
\definecolor{myblue}{HTML}{07617D}
\definecolor{mywhite}{HTML}{F5F5F5}
\definecolor{mydarkgray}{HTML}{424242}

%%%%%%%%%%%%%%%%%%%%%%%%%%%%%%%%%%%%%%%%%%%%%%%%%%%%%%%%%%%%%%%%%%%%%%%%%%%%%%%%
%% Renewed commands
%%%%%%%%%%%%%%%%%%%%%%%%%%%%%%%%%%%%%%%%%%%%%%%%%%%%%%%%%%%%%%%%%%%%%%%%%%%%%%%%
\renewcommand{\headrulewidth}{0.4pt}
\renewcommand{\footrulewidth}{0.4pt}

%%%%%%%%%%%%%%%%%%%%%%%%%%%%%%%%%%%%%%%%%%%%%%%%%%%%%%%%%%%%%%%%%%%%%%%%%%%%%%%%
%% Hyperlinks configuration
%%%%%%%%%%%%%%%%%%%%%%%%%%%%%%%%%%%%%%%%%%%%%%%%%%%%%%%%%%%%%%%%%%%%%%%%%%%%%%%%
\hypersetup{
    bookmarks=true,         % show bookmarks bar?
    unicode=false,          % non-Latin characters in Acrobat’s bookmarks
    pdftoolbar=true,        % show Acrobat’s toolbar?
    pdfmenubar=true,        % show Acrobat’s menu?
    pdffitwindow=false,     % window fit to page when opened
    pdfstartview={FitH},    % fits the width of the page to the window
    pdftitle={My title},    % title
    pdfauthor={Author},     % author
    pdfsubject={Subject},   % subject of the document
    pdfcreator={Creator},   % creator of the document
    pdfproducer={Producer}, % producer of the document
    pdfkeywords={keyword1, key2, key3}, % list of keywords
    pdfnewwindow=true,      % links in new PDF window
    colorlinks=true,        % false: boxed links; true: colored links
    linkcolor=myblue,        % color of internal links (change box color with linkbordercolor)
    citecolor=mygreen,      % color of links to bibliography
    filecolor=myred,        % color of file links
    urlcolor=myblue         % color of external links
}

%%%%%%%%%%%%%%%%%%%%%%%%%%%%%%%%%%%%%%%%%%%%%%%%%%%%%%%%%%%%%%%%%%%%%%%%%%%%%%%%
%% Listings configuration
%%%%%%%%%%%%%%%%%%%%%%%%%%%%%%%%%%%%%%%%%%%%%%%%%%%%%%%%%%%%%%%%%%%%%%%%%%%%%%%%
\lstset{%
    backgroundcolor=\color{mywhite},   % choose the background color; you must add \usepackage{color} or \usepackage{xcolor}; should come as last argument
    basicstyle=\ttfamily,            % the size of the fonts that are used for the code
    breakatwhitespace=false,         % sets if automatic breaks should only happen at whitespace
    breaklines=true,                 % sets automatic line breaking
    captionpos=b,                    % sets the caption-position to bottom
    commentstyle=\color{mygreen},    % comment style
    deletekeywords={...},            % if you want to delete keywords from the given language
    escapeinside={\%*}{*)},          % if you want to add LaTeX within your code
    extendedchars=true,              % lets you use non-ASCII characters; for 8-bits encodings only, does not work with UTF-8
    frame=l,	                     % adds a frame around the code
    framerule=1pt,                   % controls the width of the rule
    keepspaces=true,                 % keeps spaces in text, useful for keeping indentation of code (possibly needs columns=flexible)
    keywordstyle=\color{myblue},       % keyword style
    language=Prolog,                 % the language of the code
    morekeywords={*, use_module, append},                % if you want to add more keywords to the set
    numbers=left,                    % where to put the line-numbers; possible values are (none, left, right)
    numbersep=8pt,                   % how far the line-numbers are from the code
    numberstyle=\footnotesize\color{mydarkgray}, % the style that is used for the line-numbers
    rulecolor=\color{mydarkgreen},         % if not set, the frame-color may be changed on line-breaks within not-black text (e.g. comments (green here))
    showspaces=false,                % show spaces everywhere adding particular underscores; it overrides 'showstringspaces'
    showstringspaces=false,          % underline spaces within strings only
    showtabs=false,                  % show tabs within strings adding particular underscores
    stepnumber=1,                    % the step between two line-numbers. If it's 1, each line will be numbered
    stringstyle=\color{mymauve},     % string literal style
    tabsize=2,	                     % sets default tabsize to 2 spaces
    title=\lstname,                  % show the filename of files included with \lstinputlisting; also try caption instead of title
    xleftmargin=\parindent           % sets the position relative to left margin
}
%%%%%%%%%%%%%%%%%%%%%%%%%%%%%%%%%%%%%%%%%%%%%%%%%%%%%%%%%%%%%%%%%%%%%%%%%%%%%%%%
%% Start of document
%%%%%%%%%%%%%%%%%%%%%%%%%%%%%%%%%%%%%%%%%%%%%%%%%%%%%%%%%%%%%%%%%%%%%%%%%%%%%%%%
\begin{document}

%%%%%%%%%%%%%%%%%%%%%%%%%%%%%%%%%%%%%%%%%%%%%%%%%%%%%%%%%%%%%%%%%%%%%%%%%%%%%%%%
%% Cover page
%%%%%%%%%%%%%%%%%%%%%%%%%%%%%%%%%%%%%%%%%%%%%%%%%%%%%%%%%%%%%%%%%%%%%%%%%%%%%%%%
\pagenumbering{gobble}
\thispagestyle{empty}

\begin{titlepage}
    \begin{figure}
        \begin{minipage}{0.3\textwidth}
            \begin{figure}[H]
                \includegraphics[scale=0.5]{BFH_Logo}
            \end{figure}
        \end{minipage}
        \hfill
        \begin{minipage}[t]{0.5\textwidth}
            \begin{flushright}
                Berner Fachhochschule\\
                Technik und Informatik\\
                \bigskip
                Herbstsemester 16/17\\
                \bigskip
                Modul: BTI7311 - Informatik Seminar\\
                Betreuer: Kai Brünnler\\
            \end{flushright}
        \end{minipage}
    \end{figure}
    \setlength{\textfloatsep}{4cm}
    {\centering
    {\scshape\large Seminararbeit \par}
    {\scshape\LARGE\thema\par}}
    \vfill
    \begin{figure}[!hb]
        \begin{minipage}[t][][b]{0.4\textwidth}
            von\\
            Carlos Arauz\\
            5. Semester Informatik\\
            Bärenfelserstrasse 9\\
            4057 Basel\\
            arauc1@bfh.ch\\
            Tel.: +41 78 920 41 89
        \end{minipage}
        \begin{minipage}[t][][b]{0.4\textwidth}
            {\color{white}von}\\
            Angelo Campanile\\
            5. Semester Informatik\\
            Bottigenstrasse 2A\\
            3018 Bern\\
            campa1@bfh.ch\\
            Tel.: +41 79 696 09 16
        \end{minipage}
    \end{figure}
\end{titlepage}
\newpage

\pagenumbering{roman}       % set the page number to roman numbers
\setcounter{page}{2}        % start the page numbering from 2

%%%%%%%%%%%%%%%%%%%%%%%%%%%%%%%%%%%%%%%%%%%%%%%%%%%%%%%%%%%%%%%%%%%%%%%%%%%%%%%%
%% Preface
%%%%%%%%%%%%%%%%%%%%%%%%%%%%%%%%%%%%%%%%%%%%%%%%%%%%%%%%%%%%%%%%%%%%%%%%%%%%%%%%
\noindent
{\scshape\LARGE Vorworte \par}
\vspace{10mm}
\noindent
\justify
Diese Seminararbeit ist eine kurze Beschreibung, die in die Welt des Bitcoins führt.
Es wird angenommen, dass der Leser zumindest schon mal das Wort Bitcoin gehört hat,
sich aber nie vertieft damit befasst hat.

\noindent
Diese Arbeit soll dem Leser ein grösseres Verständnis von Bitcoin geben und ihm auch die
Vor- und Nachteile dieses Systems bekannt geben.

%\cite{bibo1}
%\begin{thebibliography}{9}
%\bibitem{bibo1}
%Michel Goossens, Frank Mittelbach, and Alexander Samarin.
%\textit{The \LaTeX\ Companion}.
%Addison-Wesley, Reading, Massachusetts, 1993.
%\end{thebibliography}
\newpage

%%%%%%%%%%%%%%%%%%%%%%%%%%%%%%%%%%%%%%%%%%%%%%%%%%%%%%%%%%%%%%%%%%%%%%%%%%%%%%%%
%% Index
%%%%%%%%%%%%%%%%%%%%%%%%%%%%%%%%%%%%%%%%%%%%%%%%%%%%%%%%%%%%%%%%%%%%%%%%%%%%%%%%
\tableofcontents
\newpage

% \listoffigures
% \newpage

%%%%%%%%%%%%%%%%%%%%%%%%%%%%%%%%%%%%%%%%%%%%%%%%%%%%%%%%%%%%%%%%%%%%%%%%%%%%%%%%
%% Header and Footer
%%%%%%%%%%%%%%%%%%%%%%%%%%%%%%%%%%%%%%%%%%%%%%%%%%%%%%%%%%%%%%%%%%%%%%%%%%%%%%%%
\pagestyle{fancy}
\fancyhead[L]{\bfh}
\fancyhead[R]{\modul{Informatik Seminar}}
\fancyfoot[L]{\thema}
\fancyfoot[R]{\semester{2016/2017}}

\pagenumbering{arabic}  % sets the page numbers to arabic and restarts it

%%%%%%%%%%%%%%%%%%%%%%%%%%%%%%%%%%%%%%%%%%%%%%%%%%%%%%%%%%%%%%%%%%%%%%%%%%%%%%%%
%% Capitel 000
%%%%%%%%%%%%%%%%%%%%%%%%%%%%%%%%%%%%%%%%%%%%%%%%%%%%%%%%%%%%%%%%%%%%%%%%%%%%%%%%

\vspace*{1mm}
\section*{Einleitung}
\addcontentsline{toc}{section}{Einleitung}
\vspace{-10mm}
\noindent\rule{0.8\textwidth}{0.4pt}

\vspace{5mm}

\noindent
In der letzten Dekade hat sich der Begriff Kryptowährung schnell verbreitet. Unter Kryptowährung
(auch als Kryptogeld bezeichnet) versteht man digitales Geld, welches durch die Anwendung
der Prinzipien der Kryptographie ein verteiltes, dezentrales und sicheres digitales Zahlungssystem
verwirklicht.

\noindent
Seit der Erscheinung des ersten öffentlich gehandelten Kryptogeldes, des Bitcoins, wurde eine grosse
Menge weiterer Kryptowährungen implementiert und über 3000 von denen sind zurzeit in Verwendung. In
den folgenden Kapiteln werden wir uns mit dem seit 2009 gehandelten Bitcoin beschäftigen.

\noindent
Aber damit die ersten geh versuche mit Bitcoin gelingen, muss man sich zuerst die Frage beantworten
können \emph{\dq Was ist Bitcoin eigentlich?\dq}


\newpage


%%%%%%%%%%%%%%%%%%%%%%%%%%%%%%%%%%%%%%%%%%%%%%%%%%%%%%%%%%%%%%%%%%%%%%%%%%%%%%%%
%% Capitel 001
%%%%%%%%%%%%%%%%%%%%%%%%%%%%%%%%%%%%%%%%%%%%%%%%%%%%%%%%%%%%%%%%%%%%%%%%%%%%%%%%

\section*{Was ist Bitcoin?}
\addcontentsline{toc}{section}{Was ist Bitcoin?}
\vspace{-10mm}
\noindent\rule{0.8\textwidth}{0.4pt}

\vspace{5mm}
\noindent
Bitcoin ist nichts anderes als ein Computerprogramm, das von Satoshi Nakamoto entwickelt wurde und an dem 3. Januar 2009 ins Leben gerufen wurde.
Es ist eine digitale Währung die noch sehr jung ist, noch in den Kinderschuhen steckt und laufend weiterentwickelt wird.
Es beruht auf einem dezentralen Peer-to-Peer Zahlungssystem, was bedeutet, dass keine zentrale stelle benötigt wird um die Geldbeträge zu überweisen.
Natürlich gibt es eine Kontrolle wie Transaktionen kontrolliert werden können, doch dafür wird die in den Code einprogrammierte Kryptografie verwendet.
Bitcoin wird nicht durch eine Person kontrolliert, sondern durch alle Nutzer auf der Welt.

\noindent
Bitcoin ist nur als Programm vorhanden und nicht in irgend etwas materiellem. Gleichwohl konnte sich der Wert von einem Bitcoin Jahr für Jahr vergrössern.
Wieso aber halten wir den Bitcoin für etwas Wertvolles obwohl es ihn in der realen Welt gar nicht gibt? Nun die Menschen geben Sachen eine Wert, wenn
man glaubt und das Vertrauen hat, dass etwas einen bestimmten Wert hat und andere das auch so sehen. Da man also angefangen hat mit Bitcoins zahlen zu können,
hat es somit eine Wert erhalten, da man sich somit verschiedene Sachen kaufen kann. Und solange es Menschen gibt, die den Bitcoin für Wertvoll halten, wird sich
an dem dem System nicht viel ändern.

\noindent
Da Bitcoin ein Open-Source Projekt ist, kann also jeder den Code des Systems einsehen. Dies hat den Vorteil, dass Fehler oder Sicherheitslücken schnell
gefunden und behoben werden. Somit kann auch jeder die Software einsehen, sie aber nicht den anderen erzwingen, da jeder selbst eintscheiden kann welche
Version sie benutzen wollen. Es ist jedoch nötig, dass alle Nutzer die gleiche Software brauchen, die die gleichen Regeln enthält. Nur so kann man mit
Bitcoins handeln. Und da das System nur dann funktioniert, wenn man sich einig ist, ist es natürlich von grossem interesse, dass das so bleibt.

\noindent
Um Bitcoin zu benutzen, braucht man sich nur das Bitcoin-Client Programm herunterzuladen und schon kann man sich Bitcoins kaufen. Das System wird von
vielen geschätzt und ist für viele Personen interessant, da man, um mit Bitcoins zu handeln nur einen Computer und eine Internet verbindung benötigt.
Somit sind Perosnen, welche in einem Gebiet leben, wo man keine Bank in der nähe hat, nicht von dem System ausgeschlossen.

\noindent
Das System setzt einen grossen Wert auf anonymität, was bedeutet, dass sich jeder ein \emph{\dq Bitcoin-Konto \dq} erstellen kann, da es ja nur ein
Programm ist, das man sich herunterladen muss. Deshalb werden Bitcoins natürlich auch für Illegale Zwecke verwendet, wie z.B. im Darknet. Zudem ist eine
Geldüberweisung nicht mehr umkehrbar, wie z.B. bei einer Kredirkarten buchung.

\noindent
Damit nicht eine Inflation entstehen kann wie bei einer realen Währung, es also zu viel Geld im umlauf hat. Das System beugt dies vor indem es die gesammte Menge
an exstierenden Bitcoins schon vordefiniert hat auf 21 Millionen. Die Gesamtzahl der derzeit existierenden Bitcoins beträgt ca. 16 Millionen. Und jeden Tag werden
neue generiert, bis diese 21 Millionen erreicht wurden.

\newpage


%%%%%%%%%%%%%%%%%%%%%%%%%%%%%%%%%%%%%%%%%%%%%%%%%%%%%%%%%%%%%%%%%%%%%%%%%%%%%%%%
%% Capitel 002
%%%%%%%%%%%%%%%%%%%%%%%%%%%%%%%%%%%%%%%%%%%%%%%%%%%%%%%%%%%%%%%%%%%%%%%%%%%%%%%%

\vspace*{1mm}
\section*{Wie funktioniert Bitcoin?}
\addcontentsline{toc}{section}{Wie funktioniert Bitcoin?}
\vspace{-10mm}
\noindent\rule{0.8\textwidth}{0.4pt}

\vspace{5mm}

\vspace*{1mm}
\subsection*{Hash-Funktionen}
\addcontentsline{toc}{subsection}{Hash-Funktionen}
\vspace{-10mm}
\noindent\rule{0.8\textwidth}{0.4pt}

\vspace{5mm}





\newpage

\vspace*{1mm}
\subsection*{Transaktion}
\addcontentsline{toc}{subsection}{Transaktion}
\vspace{-10mm}
\noindent\rule{0.8\textwidth}{0.4pt}

\vspace{5mm}

\noindent
In diesem Abschnitt versuchen wir zu erklären, wie genau eine solche Transaktion, also eine Geldüberweisung, Schritt für Schritt durchgeführt wird.
Dies zeigen wir anhand eines Beispiels bei dem ein Benutzer einem anderen Benutzer eine Menge von Bitcoins sendet. Diese Benutzer nennen wir zur Verständigung
\emph{\dq Sender \dq} und \emph{\dq Empfänger \dq}.

\noindent
Damit man überhaupt Bitcoins transferieren können, brauchen wir den Bitcoin client zu installieren. Dieser Client kann man sich auf dieser Seite herunterladen: \emph{\dq https://bitcoin.org/de/ \dq}.
Durch herunterlade der Software hat man sich jetzt eine Wallet erstellt, also ein Portemonnaie für die Bitcoins. Leider kann zu diesem Zeitpunkt noch nicht mit Bitcoins handeln.
Dafür muss man sich nämlich eine Adresse generieren lassen und diese ist einmalig. Mit dieser Adresse kann man nun Bitcoins transferieren.
Man kann sich beliebig viele Adressen generieren lassen, die zur gleichen wallet gehören. Es wird sogar empfohlen für jede Transaktion eine neue Adresse zu gebrauchen zum Schutz der Anonymität.
Beim Erstellen dieser Adresse wird automatisch ein Private- und ein Public-key generiert. Der Private-key ist der einzige Zugang zu der Adresse, sollte man den also verlieren, hat man keine Möglichkeit mehr
auf die Bitcoins zuzugreifen. Der Private key ist auch stark mit der Adresse gekoppelt, also mit dem Private-key kann man die Adresse berechnen aber nicht umgekehrt. Das bedeutet sollte jemand fremdes in den
Besitz dieses Private-keys gelangen, so findet er auch die Adresse und hat somit Zugang zu all sich darin befindenden Bitcoins. Es gibt keine Möglichkeit Bitcoins ohne Private-key zurück zu erhalten.
Diese Bitcoins bleiben zwar im System erhalten, jedoch können sie nicht mehr abgerufen werden. Der Public-key wird dazu verwendet, um zu validieren, ob die Transaktion wirklich von diesem Benutzer
erstellt wurde. Dies wird in einem späteren Kapitel genauer erläutert.

\noindent
Da in unserer Wallet noch keine Bitcoins vorhanden sind, können wir auch keine transferieren. Es gibt verschiedene Möglichkeiten an Bitcoins zu gelangen.
Man kann z.B. Bitcoins als Zahlung von jemandem akzeptieren und somit Bitcoins erhalten oder man geht auf einer Bitcoin-Webseite, wo man mit realem Geld Bitcoins kaufen kann.
Es gibt auch die Möglichkeit, dass man Bitcoins durch sogenanntes \emph{\dq Minen \dq} erhält, das wird in einem späteren Kapitel noch behandelt.

\noindent
Auf dem GUI des Bitcoin Programms kann der Sender nun die Adresse des Empfängers angeben und den Betrag der Bitcoins, der übertragen werden soll.
Sobald er auf Senden drückt, wird die Transaktion dem System gesendet. Uns interessiert aber was hinter dem GUI passiert und wie sich diese Transaktion zusammenstellt.
Eine Transaktion besteht aus verschiedenen Elementen. Nämlich aus einer Message, einer Signature und einem Public-key. Diese werden nun einzeln genauer erläutert:

\noindent
Die Message:
Die Message besteht selber aus verschiedenen Elementen die alle Informationen enthalten. Damit es möglich ist dass der Sender dem Empfänger eine Gewisse Anzahl Bitcoins senden kann, muss natürlich
der Sender diese Bitcoins auch haben. Nehmen wir an, der Sender, der 100 Bitcoins besitzt, möchte dem Empfänger 90 Bitcoins senden, dann wird zuerst geschaut von wo der Sender diese Bitcoins eigentlich hat. Dies ist möglich da
alle Transaktionen die jemals gemacht wurden in der Blockchain gespeichert wurden. Also sehen wir dort z.B. drei Transaktionen mit je 30 Bitcoins. diese drei Transaktionen werden nun je als Hash in der
Message gespeichert und werden als Input dieser Message angeschaut. Als Output werden auch wieder Transaktionen als Hashwert eingefügt. Dazu zählt die Information dass der Sender dem Empfänger 90 Bitcoins senden
will. Diese Information braucht die Anzahl der zu sendenden Bitcoins und den Public-key des Empfängers. Als zweiten Output Wert dieser Message wird folgendes gebraucht, die Information wie viele Bitcoins wieder
zurück zum Sender gehen. Da der Sender ja 100 Bitcoins besitzt aber nur 90 senden will, wird noch angegeben, wie viele Bitcoins wieder zurück zum Sender gelangen. Dazu braucht man den Public-key des Senders,
und die Menge der zurückgehenden Bitcoins. In diesem Fall werden wir annehmen, dass das 9 Bitcoins sein werden. Die Frage, die man sich jetzt sofort stellt ist, was passiert mit dem letzten verbleibenden
Bitcoin? Nun dieser wird als Transaktionsgebühr verwendet und geht an den Benutzer, der die Transaktion löst, aber dazu später mehr. Nun haben wir alle Elemente der Message. Und der Hash dieser Message, ist
zugleich der Name der Transaktion, wie auch der Wert der weiterverwendet wird.

\noindent
Die Signature:
Die Signature ist genau gleich wie eine Unterschrift in der realen Welt. Mit der Unterschrift bestätigt man also, dass irgendein Dokument von sich selber und niemand anderes erstellt wurde und gültig ist.
Bei der Signature beim Bitcoin ist das genau das gleiche. Man signiert eine Transaktion, um deren Wahrheitswert zu belegen und zu definieren, dass diese Transaktion zu einem selber gehört. Die Signature
setzt sich aus dem Hash der Message und dem Private-key zusammen. Diese zwei Elemente generieren einen weiteren Hash, den wir die Signature nennen. Diese Signature ist nun genau dieser Transaktion
und dieser Message gebunden. Also sie ist nicht nur dem Sender zugeteilt, was ja eigentlich logisch sein muss, da der Private-key des Senders gebraucht wird, sondern auch der Message selber. Denn wird
dem Private-key eine andere Message hinzugefügt, so ändert sich auch die Signature.

\noindent
Der Public-key
Der Public-key wird nur verwendet um zu verifizieren ob die Message und die Signature wirklich zueinander gehören und somit die Transaktion verifiziert ist.

\noindent
Da die Transaktion nun erstellt wurde kommen wir nun zum Teil, wo die Miner ins Spiel kommen.

\vspace*{1mm}
\subsubsection*{Mining}
\addcontentsline{toc}{subsubsection}{Mining}
\vspace{-10mm}
\noindent\rule{0.8\textwidth}{0.4pt}

\vspace{5mm}

\noindent
Mining wird in diesem System gebraucht, um sicherzustellen, dass alle Nutzer synchronisiert sind. Dazu werden alle in den letzten 10 Minuten erstellten Transaktionen in einen Transaktionsblock gelegt.
Nun werden die gesammelten Daten in einen Kryptographischen Hash gewandelt. Dieser Hash beinhaltet Informationen über den letzten gelösten Block, sowie eine Zufallsgenerierte Zahl, die
\emph{\dq Nounce \dq} genannt wird. Nun haben die Miner die Aufgabe, die Nounce herauszufinden, die diese Hashvalue verursacht. Das System ist so programmiert dass die Nounce eine Value
generieren muss, die mit einer bestimmten Anzahl nullen beginnen muss. Das Ratsel ist auf die Schwierigkeit eingestellt, dass man etwa 10minuten für einen Transaktionblock braucht um die Lösung
zu finden. Es ist klar, dass immer neue technologien entwickelt werden und schnellere Computer existieren werden, deshalb werden somit die Rätsel immer schneller gelöst werden. Um dieses Problem vorzubeugen,
ändert man einfach die Schwierigkeit der Nounce, damit der vorgang immer etwa 10 Minuten dauert. Der Miner, der das Rätsel gelöst hat, erhält als Belohnung ein paar Bitcoins.

\noindent
Früher brauchten die Miner Ihren Prozessor um die Rätsel zu lösen, danach hatte man gemerkt, dass mit Graphikkarten die Rätsel viel schneller gelöst werden konnte und das ging so weiter,dass
Heutzutage sogar spezialisierte Hardware existiert, die extra dafür entwickelt wurden, Bitcoins zu minen. Dadurch wurde es immer schwieriger als Individuum durch Minen Bitcoins zu erhalten.
Deshalb wurden sogenannte Mining-Pools erschaffen, wo viele Miner zusammen versuchen das Rätsel zu lösen und somit die Wahrscheinlichkeit höher ist dass man das Rätsel löst. Die Belohnung wird dann
gleichermassen an alle Miner im Pool aufgeteilt.

\noindent
Das Mining trägt dazu bei, die Sicherheit im System aufrecht zu erhalten, indem es durch die Lotterie es verhindert, das irgend jemand einen Transaction Block in die Blockchain hinzufügen kann.
Da Jeder Block in der Blockchain informationen aus dem vorherigen Block enthält, sichert sich das System vor änderungen ab. Sollte also jemand versuchen einen Block in der Blockchain
zu ändern, dann würde das sofort auffallen und als ungültig erklärt werden, da der Hashwert nicht stimmen würde. Zudem verhindert es, dass Transaktionsblöcke rückgängig gemacht werden, da man
sonst alle Transaktionen in diesem Block neu schreiben müsste.

\noindent
Also ein Miner verifiziert, ob eine Transaktionüberhaupt gültig ist indem er den Public-key mit der Message und der Signature vergleicht. Trifft dies zu, so wird die Transaktion in einen Block gespeichert,
wo alle Transaktionen gespeichert werden, welche in den Letzten 10 Minuten aufgegeben wurden.

\noindent
Die Transaktion wird in ein Kryptographisches Hash Rätsel umgewandelt. Dabei werden Informationen aus dem letzten block in der Blockchain gebraucht, den Informationen der Transaction, und einer Nounce.
Eine Nounce ist eine Zufallsgenerierte Zahl, welche zum Hash Rätsel hinzugefügt wird, dass dieser so Zufällig wie möglich ist. Dieses rätsel generiert nun ein Resultat dem sogenannten \emph{\dq Target \dq}.
Das Target ist so aufgebaut, dass die Lösung z.B. mit einer gewissen Anzahl Nullen beginnen muss. Die Miner wählen nun die Transaktion aus, welche sie lösen möchten und versuchen einen Wert, nennen wir ihn
\emph{\dq Key \dq}, herauszufinden, welches es schafft, dass die Lösung kleiner als der Target ist.

\noindent
Findet nun jemand die Lösung, also den Key, der das Rätsel löst, so werden alle Transaktionen dieses Blocks zu der Blockchain hinzugefügt. So erhält dieser Miner die Gewinnsumme der Bitcoins und Broadcastet
seine Lösung allen anderen Nutzer im Netzwerk. Diese Verifizieren, ob die Lösung korrekt ist und fügen diese in Ihre Blockchain hinzu.

\noindent
Wie wir nun verhindert, dass jemand mit einem extrem starken Rechner, jede Transaktion lösen kann, und somit eigentlich sleber bestimmen kann, welche Transaktion als nächste genommen werden muss.
Nun dies wird so gemacht, dass jeder Versuch diesen Key zu finden als ein \emph{\dq Vote \dq} angesehen wird, und so kann man verhindern, dass jemand 1000 mal mehr votet als andere.



\newpage




\newpage

\vspace*{1mm}
\subsection*{Block-Chain}
\addcontentsline{toc}{subsection}{Block-Chain}
\vspace{-10mm}
\noindent\rule{0.8\textwidth}{0.4pt}

\vspace{5mm}

\noindent
Ein weiterer wichtiger Term ohne das, die Bitcoin Währung nicht funktionieren würde, ist die Blockchain.
Die Blockchain, also auf Deutsch \emph{\dq Block - Kette\dq}, ist eine Auflistung aller Transaktionen, die jemals im Bitcoin System ausgeführt wurde.
Die momentane grösse der Blockchain beträgt momentan um die 100 GB und wird von Tag zu Tag grösser.

\noindent
Wenn wir in der realen Welt eine Geldüberweisung machen wollen, müssen wir das über einen Drittanbieter wie z.B. einer Bank erledigen. Mit der Blockchain
versucht man dies zu umgehen. Zudem erreicht man mit der Blockchain, dass die Transaktionen viel schneller überwiesen werden als mittels einer Bank z.B. Zudem
spart man sich viel Geld, da die Gebühren im Bitcoin System viel niedriger sind als bei einer Bank.

\noindent
Also Die Blockchain ist eine Liste von zusammenhängenden Transaktionen, welche nicht an  einer Zentralen Stelle gespeichert sind, da das System ja dezentral halten möchte,
von der jeder Benutzer im Netzwerk eine Kopie davon hat.

\noindent
Noch ein Vorteil der Blockchain ist, dass diese nicht manipuliert werden kann, da jeder Block Informationen des vorgehenden Block enthält. Sollte also jemand versuchen einen Block in der Blockchain
zu ändern, so würde sich auch der Hash-wert ändern, was dazu führt, dass dieser neue Block als für nicht gültig erklärt wird.

\newpage



%%%%%%%%%%%%%%%%%%%%%%%%%%%%%%%%%%%%%%%%%%%%%%%%%%%%%%%%%%%%%%%%%%%%%%%%%%%%%%%%
%% Capitel 003
%%%%%%%%%%%%%%%%%%%%%%%%%%%%%%%%%%%%%%%%%%%%%%%%%%%%%%%%%%%%%%%%%%%%%%%%%%%%%%%%

\vspace*{1mm}
\section*{Sicherheit}
\addcontentsline{toc}{section}{Sicherheit}
\vspace{-10mm}
\noindent\rule{0.8\textwidth}{0.4pt}

\vspace{5mm}





\newpage


%%%%%%%%%%%%%%%%%%%%%%%%%%%%%%%%%%%%%%%%%%%%%%%%%%%%%%%%%%%%%%%%%%%%%%%%%%%%%%%%
%% Capitel 004
%%%%%%%%%%%%%%%%%%%%%%%%%%%%%%%%%%%%%%%%%%%%%%%%%%%%%%%%%%%%%%%%%%%%%%%%%%%%%%%%

\section*{Ökonomie und Gesetz}
\addcontentsline{toc}{section}{Ökonomie und Gesetz}
\vspace{-10mm}
\noindent\rule{0.8\textwidth}{0.4pt}

\vspace{5mm}

\noindent
Wie bereits erwähnt wurde, werden neue Bitcoins durch ein in Konkurrenz stehendes und
dezentrales Verfahren namens \emph{Mining} erzeugt. Sie werden mit einer abnehmenden
und vorhersagbaren Rate generiert, dies führt dazu, dass die Anzahl neuer Bitcoins,
welche jedes Jahr produziert werden, automatisch über die Zeit halbiert, bis die Gesamtsumme
aller Bitcoins 21 Millionen erreicht.

\noindent
Der Wert des Bitcoins variiert über die Zeit. 2010 war der Wert eines Bitcoins weniger
als \$0.01 und im November 2013 erreichte der Bitcoin den bisher maximalen Wert
von \$1'250 : 1 BTC. Dass Bitcoins lediglich von einem mathematischen Algorithmus
generiert werden, lässt die Frage entstehen \emph{\dq wieso haben Bitcoins
überhaupt einen Wert?\dq}.

\noindent
Eine Währung (zum Beispiel der Dollar) beruht auf physikalischen Eigenschaften oder dem Vertauen
auf zentrale Autoritäten. Der Bitcoin hingegen basiert auf Eigenschaften der Mathematik und
er verfügt über die Charakteristik vom Geld (unter anderem Übertragbarkeit, Portabilität und
Erkennbarkeit). Dies gibt Bitcoins ihren Wert, weil sie auch in Form vom Geld genutzt werden
können und daher und wie bei allen anderen Währungen bildet sich der Wert des Bitcoins direkt und
nur durch die Menschen, die ihn als Bezahlung akzeptieren.

\noindent
Was aber der Preis eines Bitcoins bestimmt, ist, Angebot und Nachfrage. Wenn die Nachfrage nach
Bitcoins steigt, steigt der Preis, wenn sie sinkt, sinkt ebenfalls der Preis. Dass Bitcoin zur Zeit
noch ein relativ kleiner Markt ist, führt dazu, dass keine grösseren Geldbeträge nötig sind, um den
Marktpreis zu steigern oder zu verringern. Dies macht den Preis eines Bitcoins sehr volatil.

\noindent
Das Bitcoin-Protokoll kann nicht modifiziert werden ohne das Zusammenwirken fast all seiner Benutzer,
denn sie wählen, welche Software sie verwenden. Dies macht den Versuch, einer lokalen Autorität
spezielle Rechte zuzuteilen, praktisch unanwendbar. Wenn keine Autorität Bitcoins reguliert werden kann,
kann die Frage gestellt werden, ob sie wirklich legal sind. Bisher sind Bitcoins in keinem Gerichtsstand
illegal, in manchen Ländern (wie Argentinien und Russland) werden Fremdwährungen beschränkt oder
komplett verboten.

\noindent
Ein weiterer Aspekt der Bitcoins liegt hauptsächlich darin, dass Bitcoin Geld ist und Geld wurde
schon immer für legale und illegale Zwecke benutzt. Da Bitcoins für private und unumkehrbare Zahlungen
verwendet werden kann, wurden Vorwürfe gemacht, weil das Bitcoins attraktiver für Krimineller machen
könnte.


%%%%%%%%%%%%%%%%%%%%%%%%%%%%%%%%%%%%%%%%%%%%%%%%%%%%%%%%%%%%%%%%%%%%%%%%%%%%%%%%
%% Capitel 005
%%%%%%%%%%%%%%%%%%%%%%%%%%%%%%%%%%%%%%%%%%%%%%%%%%%%%%%%%%%%%%%%%%%%%%%%%%%%%%%%

\section*{Vor- und Nachteile}
\addcontentsline{toc}{section}{Vor- und Nachteile}
\vspace{-10mm}
\noindent\rule{0.8\textwidth}{0.4pt}
\subsection*{Vorteile}
\addcontentsline{toc}{subsection}{Vorteile}
\vspace{-10mm}
\noindent\rule{0.8\textwidth}{0.4pt}
\begin{description}
  \item[Zahlungsfreiheit] Bitcoin erlaubt den Benutzern, beliebige Geldbeträge
  ohne Verzögerung zu senden oder zu empfangen und zwar überall auf der Welt und
  zu jeder Zeit. Im Gegensatz zu Banken stellt Bitcoin keine Limits, wie viel Geld
  der Benutzer am Tag oder im Monat senden kann oder wie viel Geld er in seinem
  Konto haben kann. Der Benutzer kann die 365 Tage Zahlungen realisieren bzw.
  erhalten und es gibt keine Landesgrenzen. Alles, was jemand braucht, sind Bitcoins
  und ein entsprechender Bitcoin-Client.
  \item[Niedrige Gebühren] Bitcoin-Zahlungen werden derzeit mit keiner oder einer
  sehr geringen Gebühr verarbeitet. Benutzer können Gebühren zu Transkaktionen
  hinzufügen, was dazu führt, dass die entsprechende Transaktion schneller bestätigt wird.
  \item[Weniger Risiken für Händler] Da Bitcoin-Tranksaktionen sicher, unumkehrbar sind und
  keine Daten des Käufers beinhalten, werden Händler vor Verlusten geschützt, weil
  Buchungen nicht rückgängig gemacht werden können (zum Beispiel). Anderes Beispiel
  ist, dass Händler Kunden, welche keine Kreditkarten zur Verfügung haben, erreichen
  können; ohne Kreditkarten zu arbeiten, verringert das Risiko auf Betrug.
  \item[Sicherheit und Kontrolle] Da das Bitcoin-Protokoll auf mathematischen Algorithmen
  basiert und wird von keiner bestimmten Organisation gesteuert, sind Bitcoin-Transaktionen sicher,
  diese brauchen keine persönlichen Informationen der Bitcoin-Nutzer, um durchgeführt zu werden
  und somit ist die Identität des Bitcoin-Nutzers vor Identitätsdiebstahl geschützt. Es gibt keine Bitcoin-Banken und daher
  hat der Benutzer volle Kontrolle auf sein Geld. Den Bitcoin-Nutzern stehen zwei Möglichkeiten
  zur Verfügung, ihr Geld zu schützen: \emph{Backups} und \emph{Verschlüsselung}.
  \item[Transparenz und Neutralität] Das Bitcoin-Protokoll ist sicher verschlüsselt, daher kann
  keine Einzelperson oder Organisation es kontrollieren oder manipulieren. Aus diesem Grund können die Bitcoin-Nutzer
  dem Kern von Bitcoin vertrauen, dass er völlig neutral, transparent und vorhersehbar ist.

\end{description}
\subsection*{Nachteile}
\addcontentsline{toc}{subsection}{Nachteile}
\vspace{-10mm}
\noindent\rule{0.8\textwidth}{0.4pt}
\begin{description}
  \item[Grad der Akzeptanz] Vielen Menschen sind Bitcoin noch unbekannt oder sie haben davon
  gehört, möchten aber es nicht ausprobieren, weil sie nicht verstehen, wie sie funktionieren
  oder was Bitcoins eigentlich sind. Obwohl jeden Tag immer mehr Firmen Bitcoins als Zahlungsmethode
  akzeptieren, verbleibt die Liste noch klein und um vom Netzwerkeffekt profitieren zu können, müsste
  diese Liste wachsen.
  \item[Volatilität] Da Bitcoins ihre volle Potenzial wegen einer Mangel an Akzeptanz noch nicht
  erreicht haben, hängt ihr Wert davon ab, wie gefordert Bitcoins sind. Mit anderen Worten, wenn Menschen
  viel Interesse an Bitcoins haben und bereit sind, einen gewissen Betrag für sie zu zahlen, steigt der
  Wert des Bitcoin und wenn nicht, sinkt er.
  \item[Laufende Entwicklung] Die Bitcoin Software und Bitcoins selbst sind noch relativ jung. Das heisst,
  dass sich viele Funktionen noch in der Entwicklungsphase befinden. Neue Funktionen, Werkzeuge und Dienste
  werden auch eingeführt, um Bitcoin sicherer zu machen und um zu machen, dass Bitcoin eine grössere Masse
  erreicht.

\end{description}


%%%%%%%%%%%%%%%%%%%%%%%%%%%%%%%%%%%%%%%%%%%%%%%%%%%%%%%%%%%%%%%%%%%%%%%%%%%%%%%%
%% Capitel Schlusswort
%%%%%%%%%%%%%%%%%%%%%%%%%%%%%%%%%%%%%%%%%%%%%%%%%%%%%%%%%%%%%%%%%%%%%%%%%%%%%%%%

\vspace*{1mm}
\section*{Schlussworte}
\addcontentsline{toc}{section}{Schlussworte}
\vspace{-10mm}
\noindent\rule{0.8\textwidth}{0.4pt}

Die Seminararbeit zum Thema Bitcoin hat uns einen grösseren Einblick in die Welt eiens dezentralen Rechnungsnetzwerk eingeführt.
Zu Beginn der Arbeit war das Thema ziemlich komplex und kaum hatte man sich ein paar Informationen geholt, verursachten diese
viele weitere Fragen. Der Begriff Bitcoin war uns schon bekannt, jedoch hatten wir uns nie im detail damit beschäftigt, was wir
nun dank dieses Seminars tun durften. Es war Spannend zu entdecken welche Möglichkeiten Bitcoins bieten sowie die Stärken und
Schwächen eines solchen Systems herauszufinden.

\noindent
Wie schon erwähnt wurde, befindet sich Bitcoin immernoch in der Beta-phase, und es wird sehr spannend sein die weitere Entwicklung
des Systems zu beobachten. Vorallem wird es spannend sein zu sehen wie sich das System entwickelt, wenn es für die Miner plötzlich
nicht mehr rentabel wird. Und da Bitcoins eine bestimmte Menge vorprogrammiert hat wird dies einmal der Fall sein.
Wir sind uns sicher dass eine gute Lösung gefunden wird, die die Nutzer auch akzeptieren. Im moment müssen wir uns keine Sorgen machen,
da es noch genügend Bitcoins an vorrat hat.

\noindent
Die Arbeit hat viel Zeit in Anspruch genommen, doch am Ende hat es sich gelohnt.
Wahrend dieser Arbeit haben wir viel Wissen mitnehmen konnen, was wir in Zukunft
sicherlich einmal gebrauchen konnen. Dies ist nun das Ende unserer zweiten Informatik-
Seminar Arbeit, wir hoffen das wir Ihnen einen kleinen Einblick in die Welt des Bitcoins bieten durften.
Vielen Dank an die Dozenten welche uns die Moglichkeit gegeben haben dieses Thema erarbeiten zu durfen.


\nocite{*}
\newpage
\addcontentsline{toc}{section}{Literatur}
\bibliography{Bibliografie/literatur}
\bibliographystyle{alpha}
\end{document}
